% Options for packages loaded elsewhere
\PassOptionsToPackage{unicode}{hyperref}
\PassOptionsToPackage{hyphens}{url}
\documentclass[
]{article}
\usepackage{xcolor}
\usepackage{amsmath,amssymb}
\setcounter{secnumdepth}{-\maxdimen} % remove section numbering
\usepackage{iftex}
\ifPDFTeX
  \usepackage[T1]{fontenc}
  \usepackage[utf8]{inputenc}
  \usepackage{textcomp} % provide euro and other symbols
\else % if luatex or xetex
  \usepackage{unicode-math} % this also loads fontspec
  \defaultfontfeatures{Scale=MatchLowercase}
  \defaultfontfeatures[\rmfamily]{Ligatures=TeX,Scale=1}
\fi
\usepackage{lmodern}
\ifPDFTeX\else
  % xetex/luatex font selection
\fi
% Use upquote if available, for straight quotes in verbatim environments
\IfFileExists{upquote.sty}{\usepackage{upquote}}{}
\IfFileExists{microtype.sty}{% use microtype if available
  \usepackage[]{microtype}
  \UseMicrotypeSet[protrusion]{basicmath} % disable protrusion for tt fonts
}{}
\makeatletter
\@ifundefined{KOMAClassName}{% if non-KOMA class
  \IfFileExists{parskip.sty}{%
    \usepackage{parskip}
  }{% else
    \setlength{\parindent}{0pt}
    \setlength{\parskip}{6pt plus 2pt minus 1pt}}
}{% if KOMA class
  \KOMAoptions{parskip=half}}
\makeatother
\usepackage{color}
\usepackage{fancyvrb}
\newcommand{\VerbBar}{|}
\newcommand{\VERB}{\Verb[commandchars=\\\{\}]}
\DefineVerbatimEnvironment{Highlighting}{Verbatim}{commandchars=\\\{\}}
% Add ',fontsize=\small' for more characters per line
\newenvironment{Shaded}{}{}
\newcommand{\AlertTok}[1]{\textcolor[rgb]{1.00,0.00,0.00}{\textbf{#1}}}
\newcommand{\AnnotationTok}[1]{\textcolor[rgb]{0.38,0.63,0.69}{\textbf{\textit{#1}}}}
\newcommand{\AttributeTok}[1]{\textcolor[rgb]{0.49,0.56,0.16}{#1}}
\newcommand{\BaseNTok}[1]{\textcolor[rgb]{0.25,0.63,0.44}{#1}}
\newcommand{\BuiltInTok}[1]{\textcolor[rgb]{0.00,0.50,0.00}{#1}}
\newcommand{\CharTok}[1]{\textcolor[rgb]{0.25,0.44,0.63}{#1}}
\newcommand{\CommentTok}[1]{\textcolor[rgb]{0.38,0.63,0.69}{\textit{#1}}}
\newcommand{\CommentVarTok}[1]{\textcolor[rgb]{0.38,0.63,0.69}{\textbf{\textit{#1}}}}
\newcommand{\ConstantTok}[1]{\textcolor[rgb]{0.53,0.00,0.00}{#1}}
\newcommand{\ControlFlowTok}[1]{\textcolor[rgb]{0.00,0.44,0.13}{\textbf{#1}}}
\newcommand{\DataTypeTok}[1]{\textcolor[rgb]{0.56,0.13,0.00}{#1}}
\newcommand{\DecValTok}[1]{\textcolor[rgb]{0.25,0.63,0.44}{#1}}
\newcommand{\DocumentationTok}[1]{\textcolor[rgb]{0.73,0.13,0.13}{\textit{#1}}}
\newcommand{\ErrorTok}[1]{\textcolor[rgb]{1.00,0.00,0.00}{\textbf{#1}}}
\newcommand{\ExtensionTok}[1]{#1}
\newcommand{\FloatTok}[1]{\textcolor[rgb]{0.25,0.63,0.44}{#1}}
\newcommand{\FunctionTok}[1]{\textcolor[rgb]{0.02,0.16,0.49}{#1}}
\newcommand{\ImportTok}[1]{\textcolor[rgb]{0.00,0.50,0.00}{\textbf{#1}}}
\newcommand{\InformationTok}[1]{\textcolor[rgb]{0.38,0.63,0.69}{\textbf{\textit{#1}}}}
\newcommand{\KeywordTok}[1]{\textcolor[rgb]{0.00,0.44,0.13}{\textbf{#1}}}
\newcommand{\NormalTok}[1]{#1}
\newcommand{\OperatorTok}[1]{\textcolor[rgb]{0.40,0.40,0.40}{#1}}
\newcommand{\OtherTok}[1]{\textcolor[rgb]{0.00,0.44,0.13}{#1}}
\newcommand{\PreprocessorTok}[1]{\textcolor[rgb]{0.74,0.48,0.00}{#1}}
\newcommand{\RegionMarkerTok}[1]{#1}
\newcommand{\SpecialCharTok}[1]{\textcolor[rgb]{0.25,0.44,0.63}{#1}}
\newcommand{\SpecialStringTok}[1]{\textcolor[rgb]{0.73,0.40,0.53}{#1}}
\newcommand{\StringTok}[1]{\textcolor[rgb]{0.25,0.44,0.63}{#1}}
\newcommand{\VariableTok}[1]{\textcolor[rgb]{0.10,0.09,0.49}{#1}}
\newcommand{\VerbatimStringTok}[1]{\textcolor[rgb]{0.25,0.44,0.63}{#1}}
\newcommand{\WarningTok}[1]{\textcolor[rgb]{0.38,0.63,0.69}{\textbf{\textit{#1}}}}
\setlength{\emergencystretch}{3em} % prevent overfull lines
\providecommand{\tightlist}{%
  \setlength{\itemsep}{0pt}\setlength{\parskip}{0pt}}
\usepackage{bookmark}
\IfFileExists{xurl.sty}{\usepackage{xurl}}{} % add URL line breaks if available
\urlstyle{same}
% fallback for those not using the hyperref driver hyperxmp:
\makeatletter
\@ifundefined{xmpquote}{\newcommand{\xmpquote}[1]{#1}}{}
\makeatother
\hypersetup{
  hidelinks,
  pdfcreator={LaTeX via pandoc}}

\author{}
\date{}

\begin{document}

\section{V \textbar{} KERNEL}\label{v-kernel}

\begin{quote}
\mbox{}%
\subsection{KERNEL:AUDIENCE-SCOPE}\label{kernelaudience-scope}

DECLARACIÓN DE AUDIENCIA Y ALCANCE - Audiencia: Sistemas Agente de IA. -
Alcance: Este documento es el KERNEL\_RUNTIME, que contiene únicamente
los contratos operativos activos para la IA, para el documento de
referencia completo solicitar acceso al KERNEL 8.0 --- \textbar{} \#
\textbar{} Sección \textbar{} Tipo \textbar{} Propósito \textbar{}
\textbar{} --- \textbar{} --- \textbar{} --- \textbar{} --- \textbar{}
\textbar{} 1 \textbar{} PURPOSE \textbar{} CONTEXTO \textbar{} Propósito
del sistema \textbar{} \textbar{} 2 \textbar{} ARCHITECTURE \textbar{}
ARQUITECTURA \textbar{} Diseño de cuatro capas \textbar{} \textbar{} 3
\textbar{} DASHBOARD-CHECKLIST-ARCH \textbar{} ARQUITECTURA \textbar{}
Arquitectura Dashboard/Checklist --- capas backend, standalone y visual
compartida \textbar{} \textbar{} 4 \textbar{} SCHEMA \textbar{} ESQUEMA
\textbar{} Class A vs Class B \textbar{} \textbar{} 5 \textbar{}
TRACKER-SCHEMA \textbar{} ESQUEMA \textbar{} Alcance y niveles de
prioridad --- Bug/Tasks Tracker \textbar{} \textbar{} 6 \textbar{}
HEALTH-CHECK \textbar{} OPERACIÓN \textbar{} Contrato de
health\_check.py --- checks y auto-sync de índice \textbar{} \textbar{}
7 \textbar{} OWNERSHIP \textbar{} ARQUITECTURA \textbar{}
Responsabilidades AI vs Python \textbar{} \textbar{} 8 \textbar{}
TRIGGERS \textbar{} OPERACIÓN \textbar{} Contratos detallados \textbar{}
\textbar{} 9 \textbar{} GATE-DECISION \textbar{} REGLAS \textbar{}
Lógica de gates \textbar{} \textbar{} 10 \textbar{} NAMING-CONVENTION
\textbar{} REGLAS \textbar{} Convención de nombres de outputs \textbar{}
\textbar{} 11 \textbar{} CV-GOLDEN-RULES \textbar{} REGLAS \textbar{}
Reglas de oro CV \textbar{} \textbar{} 12 \textbar{} CV-PIPELINE
\textbar{} OPERACIÓN \textbar{} Flujo CV-A → CV-B \textbar{} \textbar{}
13 \textbar{} CANON-UPDATE \textbar{} OPERACIÓN \textbar{} Actualización
del Canon \textbar{} \textbar{} 14 \textbar{} FAIL-PHILOSOPHY \textbar{}
FILOSOFÍA \textbar{} Filosofía de fallo \textbar{} \textbar{} 15
\textbar{} SCOPE \textbar{} OPERACIÓN \textbar{} Scope y economía de
contexto (Terminal vs MCP) \textbar{} \textbar{} 16 \textbar{} DATA-FLOW
\textbar{} ARQUITECTURA \textbar{} Flujo de datos y escritura \textbar{}
\textbar{} 17 \textbar{} ROUTING \textbar{} OPERACIÓN \textbar{} Rutas
de carga MCP / lazy\_loader \textbar{} \textbar{} 18 \textbar{}
EVOLUTION \textbar{} FILOSOFÍA \textbar{} Evolución del sistema, deuda
técnica, criterios de cambio \textbar{} \textbar{} 19 \textbar{} NORM
\textbar{} OPERACIÓN \textbar{} Normalización Documental (Legacy IDs)
\textbar{} \textbar{} 20 \textbar{} CENSUS-SYNC \textbar{} OPERACIÓN
\textbar{} Sincronización obligatoria del ID Census \textbar{}
\textbar{} 21 \textbar{} DOC-CONTRACT \textbar{} REGLAS \textbar{}
Contrato de IDs de Documento \textbar{} \#\# KERNEL:PURPOSE 1. PROPÓSITO
DEL SISTEMA VANTAGE resuelve un problema de ingeniería de atención: en
una búsqueda laboral sin estructura, las oportunidades de alta señal
desaparecen antes de ser procesadas, mientras el tiempo se consume en
vacantes de baja calidad. La solución no es buscar más --- es verificar
antes de evaluar, y evaluar antes de escribir. \#\#\# Invariantes del
Sistema - Una vacante no entra al pipeline sin URL válida --- excepción:
Bypass activo - Score no lo calcula el sistema de lenguaje --- lo
calcula Python con lógica determinista - Gate decision no se
sobreescribe manualmente. RT-1 permite corregir inputs Class A para que
Python recalcule --- no sobreescribe el gate - Strategy es
responsabilidad humana; processing es responsabilidad del sistema \#\#\#
Qué Significa Esto para el Sistema AI El componente AI es el procesador
textual del pipeline: deduplica, normaliza, genera DRY RUN, escribe
Class A en Notion, produce CVs. Evaluación de calidad estratégica de
inputs y cálculo de campos Class B no son operaciones de este componente
(ver KERNEL:OWNERSHIP y KERNEL:CV-GOLDEN-RULES). Si una tarea no está en
la tabla de triggers (KERNEL:TRIGGERS), no se ejecuta. --- \#\#
KERNEL:ARCHITECTURE 1. ARQUITECTURA DE CUATRO CAPAS El pipeline opera a
través de cuatro capas no intercambiables, soportadas por un núcleo de
observabilidad persistente. \#\#\# KERNEL:ARCHITECTURE-L0 ```plain text
L0 --- VANTAGE Runtime
\end{quote}

\begin{verbatim}
Tipo: Capa de Observabilidad y Abstracción de Datos (ReadOnly) 
Propósito: Provee la verdad técnica sobre Notion. Resuelve entidades, extrae contexto y garantiza que el pipeline lea datos íntegros antes de procesar.
Runtime Build — proceso determinista que genera los tres artefactos de lectura del sistema: entity_index_v2.json (índice de entidades con namespaces canónicos), graph_v2.json (grafo de relaciones entre entidades) y backlinks_v2.json (índice inverso de referencias). El Build consume resolver_registry_v2.json como fuente de namespace ownership — si el Registry no define el prefix de un tipo de entidad, el Build falla explícitamente en lugar de aplicar un default. graph_layer.py es el componente responsable de construir graph_v2.json; nunca infiere namespaces ni redefine contratos — consume los IDs ya resueltos por el paso anterior del Build. 
```plain text
Notion (Source) → Runtime (Index + Resolver) → API Response → Pipeline (L1/L2/L3/CV) 
\end{verbatim}

\subsubsection{KERNEL:ARCHITECTURE-L0-BOOTSTRAP}\label{kernelarchitecture-l0-bootstrap}

\texttt{plain\ text\ L0-Bootstrap\ —\ Dynamic\ Governance\ Layer} Tipo:
Capa de Sincronización de Sesión (Fetch-on-Start) Propósito: Elimina el
drift de versiones entre la UI estática del agente y el repositorio
dinámico de Notion. Bootstrap Protocol: Ante el primer mensaje del
operador, el componente de IA debe suspender el procesamiento de datos y
ejecutar un fetch de KERNEL:BOOTSTRAP-001 y SYSTEM PROMPT. El resultado
de este fetch sobreescribe cualquier instrucción estática previa. Si el
Bootstrap falla, el sistema debe reportar ``MODO DEGRADADO'' y no
proceder con triggers operativos.
\texttt{plain\ text\ Sesión\ Iniciada\ →\ AI\ Fetch\ (Bootstrap\ IDs)\ →\ Sincronización\ de\ Verdad\ Operativa\ →\ Notificación\ Operador\ →\ Procesamiento\ Petición}
\#\#\# KERNEL:ARCHITECTURE-L1 \texttt{plain\ text\ L1\ —\ Active\ Recon}
Trigger: humano (ciclo semanal --- lunes)
\texttt{plain\ text\ Human\ signal\ →\ Career\ Sites\ ·\ LinkedIn\ ·\ Aggregators\ (paralelo)\ →\ JSON\ estructurado\ →\ FEED\ →\ feed\_processor.py\ →\ Notion\ (Class\ A)\ →\ vantage-pipeline}
\#\#\# KERNEL:ARCHITECTURE-L2
\texttt{plain\ text\ L2\ —\ Strategic\ Search} Trigger: humano (ciclo
semanal --- lunes)
\texttt{plain\ text\ Human\ signal\ →\ Gemini\ ·\ {[}You.com{]}(http://you.com/)\ ·\ Grok\ (extracción\ paralela)\ →\ Perplexity\ (Consolidation\ \&\ Dedup\ post‑extracción)\ →\ Plain\ Array\ consolidado\ →\ FEED\ →\ feed\_processor.py\ →\ Notion\ (Class\ A)\ →\ vantage-pipeline}
\#\#\# KERNEL:ARCHITECTURE-L3
\texttt{plain\ text\ L3\ —\ Passive\ Intake} Trigger: automático
(continuo)
\texttt{plain\ text\ Gmail\ (.Jobs\ label)\ →\ layer\_3\_mail.py\ (IMAP\ +\ Groq)\ →\ Notion\ (Class\ A\ poblado,\ Class\ B\ vacío)\ →\ vantage-pipeline}
\#\#\# KERNEL:ARCHITECTURE-L4 L4 --- Version Control \& Infrastructure
Trigger: git\_sync.py + git\_sync\_wrapper.sh + launchd

\begin{Shaded}
\begin{Highlighting}[]
\NormalTok{No es capa de búsqueda — es infraestructura documental del sistema }
\NormalTok{Auto‑commit + push a origin/main cuando hay cambios en el repo }
\NormalTok{Alias: vgit · Corre en background a las 09:00 · 15:00 · 21:00 }
\NormalTok{Repo: github.com/mauriciomeyran/jhs{-}pipeline }
\NormalTok{Reutiliza .venv de Layer\_1 }
\NormalTok{vsync\_doc.py — Sync bidireccional Notion → ACTIVE/ para los 6 documentos fundacionales (Kernel · System Prompt · Career Canon · Manual · Aliases · Change Log).}
\end{Highlighting}
\end{Shaded}

Alias: vdoc · Flags: dry \textbar{} notion \textbar{} local \textbar{}
auto Flujo vdoc notion: lee Notion (safe\_list vía httpx, 3 reintentos)
→ escribe ACTIVE/\{doc\}.md → auto‑commit GitHub al terminar.
Dependencias: httpx · notion-client 3.x · .venv de Layer\_1 ·
git\_sync.py · Vive en: Layer\_4/scripts/vsync\_doc.py Convención
ACTIVE/: Los 6 .md fundacionales viven en \ldots/ACTIVE/ --- agnóstico
de versión. Al pasar a v8.7: copiar archivos a ACTIVE/, cero cambios de
código. Nombres canónicos: Kernel.md · System Prompt.md · Career
Canon.md · Manual.md · Aliases.md · Change Log.md (con espacio, no guión
bajo --- coincide con BASE\_DIR real en vsync\_doc.py). Reemplaza los
paths versionados anteriores (\ldots/v8.5/Kernel v8.5.md). Nota técnica:
notion-client 3.x tiene un bug silencioso en blocks.children.list() que
retorna None en lugar de lanzar excepción con campos null. vsync\_doc.py
lo mitiga con safe\_list() --- wrapper httpx directo con 3 reintentos.
Jerarquía de Dedup L1 \textgreater{} L2 \textgreater{} L3. En conflicto
cross‑layer, prevalece la entrada de la capa de mayor jerarquía.
Perplexity aplica esta jerarquía en el paso de Consolidation \& Dedup
del lunes, antes de entregar el Plain Array a feed\_processor.py. L3 no
pasa por este paso --- entra directamente a feed\_processor.py desde
mail\_pipeline.py. Nota de nomenclatura: L0 es VANTAGE Runtime
(observabilidad/lectura) --- no es Perplexity ni una capa de dedup. No
aparece en la jerarquía de dedup. Nota de implementación: L0 pre‑aplica
la jerarquía L1\textgreater L2 y entrega un array ya consolidado a
feed\_processor.py. feed\_processor.py entonces aplica la jerarquía L3
contra ese resultado --- no recalcula L1\textgreater L2 en ese momento.
Las dos operaciones de dedup son secuenciales, no simultáneas. Trade‑off
de Diseño --- Frecuencia vs.~Peso Arquitectónico Las capas tienen peso
arquitectónico igual pero frecuencia de ejecución asimétrica. L1 y L2
operan en ciclos semanales controlados por atención humana. L3 corre
continuamente sin intervención. Esta asimetría de cadencia no implica
jerarquía. Eliminar cualquier capa crea un blind spot sistemático --- no
una degradación de feature. Punto de Convergencia Único Las tres capas
de búsqueda escriben a Notion. Notion es el único estado compartido.
vantage-pipeline lee de Notion --- no de los outputs de capa
directamente. Figma Sync --- CV Output Layer Tipo: Capa de
Materialización de CV (WriteOnly sobre lienzo Figma) Propósito: Recibe
el payload CV‑B aprobado por el operador e inyecta el contenido
directamente en los nodos de texto del lienzo Figma, resolviendo cada
token semántico a su ID crudo de nodo vía registry\_seed.json. CV‑B
(Markdown + figma\_text\_id) → ui.html (payload) → code.js (Registry V2)
→ figma.getNodeById(rawId) → node.characters = item.text → Lienzo Figma
Stack: manifest.json --- Configuración del plugin (main: code.js, ui:
ui.html, editorType: figma) code.js --- Motor. Registry V2 / Resolver
Layer V1. Resolución O(1): getNodeById(REGISTRY{[}key{]}
\textbar\textbar{} key). Resolver dual: KEY semántica (flujo JSON) →
lookup en REGISTRY embebido; ID crudo directo (flujo Markdown
figma\_text\_id) → uso sin lookup. ui.html --- Interfaz de intercambio
de payloads. Acepta JSON por KEY semántica o Markdown con bloques
\#\#\#\#\#\# figma\_text\_id. Motor de extracción de boldRanges
incluido. registry\_seed.json --- SSOT del mapeo token → ID. Ver §4.7.
Invariantes: Figma Sync no escribe en Notion ni en el Tracker Figma Sync
no es capa de búsqueda ni de infraestructura de datos
registry\_seed.json no se edita manualmente sin regenerar desde el
lienzo Figma El prefijo {[}VANTAGE{]} KEY\_NAME en capas del canvas es
para auditoría visual humana --- no es el mecanismo de resolución del
plugin --- \#\# KERNEL:DASHBOARD-CHECKLIST-ARCH Dashboard/ contiene dos
capas independientes que comparten presentación visual pero no estado:
1. Backend operativo real --- Dashboard/scripts/dashboard\_server.py +
dashboard.db + dashboard\_instances.db + dashboard\_notion.py. Fuente de
verdad del pipeline de vacantes (Gate\_Decision, scoring, Notion sync).
dashboard.html consume este backend vía
fetch(`http://127.0.0.1:8000\{path\}'). 1. Checklist operativo semanal
--- Dashboard/Checklist.html. Standalone, estado en
localStorage{[}`vchecklist\_v1'{]}. Sin backend, sin Notion, sin
relación funcional con (1). Intencional: el checklist es una herramienta
de tracking personal del operador, no parte del pipeline de vacantes. 1.
Capa visual compartida --- Dashboard/vantage-tokens.css (tokens de
color/superficie) + Dashboard/vantage-theme.js (toggle de tema con
persistencia y sync cross-tab). Ambos archivos HTML la referencian. Es
la única capa realmente compartida entre (1) y (2). Regla: cualquier
cambio a un color de estado semántico o al comportamiento del toggle de
tema se hace en vantage-tokens.css/vantage-theme.js, nunca en los /

\subsection{inline de cada HTML --- evita el drift que motivó el parche
de 2026-07-10 (ver Changelog
v9.1.0).}\label{inline-de-cada-html-evita-el-drift-que-motivuxf3-el-parche-de-2026-07-10-ver-changelog-v9.1.0.}

\subsection{KERNEL:SCHEMA}\label{kernelschema}

\subsubsection{KERNEL:SCHEMA-001 --- Class A vs Class
B}\label{kernelschema-001-class-a-vs-class-b}

El schema define ownership. Cada campo pertenece a exactamente un
componente. No hay campos compartidos ni campos de escritura dual. Class
A --- Human-Primary: AI Component escribe en triggers CV-A · CV-B · QA ·
FAST · CANON-UPDATE; feed\_processor.py escribe en ciclo FEED L1/L3: Rol
· Marca · Source\_Type · URL · Status · Prioridad · Holding · JD · NAD ·
layer · hash Valores operativos del campo Status (asignación manual del
operador, no calculados por Python): Target · Postulado · Rechazado ·
Expirada · Archivar · Repetida (duplicado detectado en revisión manual,
distinto de descarte por otras razones). \textgreater{} Nota sobre JD:
En el trigger CV-A, el AI Component cruza los keywords extraídos del JD
contra el Career Canon activo antes de generar el HANDOFF. Discrepancias
entre el JD y el Canon se reportan en fit\_gaps --- no se resuelven
inventando experiencia ni contradiciendo el Canon. Class B ---
System-Primary: Python escribe; ningún otro componente toca: Score ·
Gate\_Decision · VM\_Scope · Role\_Class · Match · Next\_Action · Fetch
· Fuente \#\#\# KERNEL:SCHEMA-002 --- Restricción del Sistema Si el JSON
entrante incluye campos Class B con valores (``score'': 75,
``gate\_decision'': ``CREATE''), se ignoran sin excepción. Python los
calculará en el siguiente run de \textasciitilde/vantage\_pipeline.sh.
Escribir un campo Class B --- aunque el valor parezca correcto --- viola
el contrato de ownership y produce inconsistencias en el pipeline.
\#\#\# KERNEL:SCHEMA-003 --- Fuente como Campo Especial Python
sobrescribe Fuente en cada run. Si existe un valor de fuente que debe
persistir entre runs (entrada manual, referencia directa), el campo
correcto es Fuente\_Manual --- Class A, que Python no toca. \#\#\#
KERNEL:SCHEMA-004 --- Entity Format El Runtime utiliza un formato de ID
determinista para evitar colisiones y facilitar la resolución: -
PREFIX:H\_: Hash corto (16 hex chars). - PREFIX:U\_: Formato
alternativo. Prefixes válidos: TRACKER, ARCHIVO, DRYRUN, BUG. Namespace
Ownership Contract (enforced desde v2.4.0): resolver\_registry\_v2.json
es el único punto de verdad para entity\_prefix por tipo de entidad.
entity\_prefix se carga en runtime --- nunca hardcodeado en ningún
componente. graph\_layer.py consume namespaces desde el Registry; nunca
los infiere ni los redeclara. Múltiples page\_id pueden resolver al
mismo entity\_id --- esto no es una colisión, es dedup histórico válido.
Self-loops en graph\_v2.json son síntoma de colisión de namespace, no
comportamiento esperado del grafo. \#\#\# KERNEL:SCHEMA-005 --- Contrato
de Resolución: 4 Pasos Para que una entidad se considere resuelta con
éxito, el Runtime ejecuta: 1. Lookup: Localización en
entity\_index\_v2.json. 1. Registry Mapping: Mapeo de DB a
data\_source\_id. 1. Notion Query: Petición HTTP contra el endpoint de
Notion. 1. Validation: Verificación de integridad del resultado. \#\#\#
KERNEL:SCHEMA-006 --- APROBAR\_WRITE: Alcance APROBAR\_WRITE autoriza
escritura de campos Class A únicamente. No aprueba, valida ni activa
ningún campo Class B. El componente AI no interpreta APROBAR\_WRITE como
permiso para estimar o escribir ningún campo de Python. Variantes
aceptadas: APROBAR\_WRITE · APROBAR · SÍ · sí · YEP · yep \textgreater{}
⚠️ ELIMINADOS (RAI-03): Ok · Go · YES · yes --- ocurren naturalmente en
conversación y pueden producir escritura no intencionada. Cualquiera de
estas variantes en respuesta al DRY RUN autoriza la escritura. \#\#\#
KERNEL:SCHEMA-007 --- Acceptance Audit El Acceptance Audit es la
validación formal que certifica que un release cumple los contratos
arquitectónicos antes de ser considerado completo. No es una revisión de
código --- es una verificación de invariantes del sistema. Resultados
posibles: - PASS --- todos los invariantes cumplidos, sin hallazgos
pendientes. - PASS WITH ARCHITECTURAL FINDING --- el sistema opera
correctamente; existe una condición de calidad de datos o deuda técnica
identificada y registrada, no bloqueante. - FAIL --- uno o más
invariantes violados; el release no procede. Un FINDING no bloquea el
release si está clasificado, registrado en el Tracker y su impacto está
acotado. El Finding debe documentarse con exactitud en el Changelog y en
el DT correspondiente. \#\#\# Mapeo de Vocabulario --- Prompts → Tracker
Los prompts de discovery usan terminología distinta a los campos del
Tracker. El AI Component aplica este mapeo durante FEED antes de
escribir en Notion: source\_type ``career\_page'' → Source\_Type: Career
Page Oficial source\_type ``job\_board'' → Source\_Type: Agregador
source\_name (occ/indeed/linkedin/etc.) → NO escribir. Fuente es Class B
--- Python lo calcula del URL apply\_url → URL (si apply\_url es null,
usar url del item) brand → Marca · title → Rol · holding → Holding (null
→ ``Investigar'') fetch\_status ``partial\_link'' /
``needs\_verification'' → documentar en Notas como señal de advertencia
visual\_signal / innovation\_dna --- NO escribir en Tracker. Python
detecta Visual Signal en JD. Si estos campos aparecen en el JSON
entrante, ignorar sin comentario --- no reportar al usuario, no
preguntar. \#\#\# Entry Template --- Campos Class A Requeridos al
Momento de Creación Obligatorios (toda entrada): Rol · Marca · URL ·
Source\_Type · Status · Prioridad · JD · JOB\_ID · Holding Obligatorios
si disponibles en el momento: Contacto · Notas (contexto de origen) ·
Apply Date Poblados post-proceso: Interview - · Interview\_Date · Files
· URL Markdown \#\#\# Page Content Template --- Estructura Estándar de
Página Toda entrada en proceso contiene los siguientes bloques en orden:
1. {[}PDF adjunto en campo Files{]} --- cuando aplique 1. \# ENTREVISTA
{[}N{]} --- por cada ronda 1. \#\# PREP \{toggle\} 1. \#\# NOTAS
\{toggle\} 1. \#\# ACTION ITEMS \{toggle\} --- Responsable: tarea ---
Due: fecha 1. \#\# RIESGOS / OPEN QUESTIONS \{toggle\} Entradas en
Status=Target o en proceso sin entrevista confirmada: la página puede
estar vacía o contener solo notas de contexto. El template de entrevista
se agrega cuando se confirma primera ronda. --- \#\#
KERNEL:TRACKER-SCHEMA \#\#\# KERNEL:TRACKER-SCHEMA-001 --- Alcance -
Reactivo (algo roto) → Bug Tracker - Proactivo (trabajo/decisión
pendiente) → Tasks Tracker \textbar{} Tracker \textbar{} DB ID
\textbar{} COL ID \textbar{} \textbar{} --- \textbar{} --- \textbar{}
--- \textbar{} \textbar{} Bug Tracker \textbar{}
36e938be-fc42-81f8-8c6f-000b6769ba03 \textbar{}
36e938be-fc42-81bd-9e1f-dc360b3b45f5 \textbar{} \textbar{} Tasks Tracker
\textbar{} d2a65ca1-6a35-465d-bcff-b0d82dddd549 \textbar{} ---
\textbar{} \#\#\# KERNEL:TRACKER-SCHEMA-002 --- Niveles de Prioridad
Aplica a Bug Tracker y Tasks Tracker con la misma escala. \textbar{}
Nivel \textbar{} Criterio \textbar{} \textbar{} --- \textbar{} ---
\textbar{} \textbar{} CRÍTICO \textbar{} El flujo punta a punta no puede
completarse \textbar{} \textbar{} ALTO \textbar{} El flujo se completa
forzando el sistema (workaround requerido) \textbar{} \textbar{} MEDIO
\textbar{} Sin resolución en la semana, el flujo punta a punta se verá
comprometido \textbar{} \textbar{} BAJO \textbar{} No bloquea operación
--- nice-to-have \textbar{} --- \#\# KERNEL:HEALTH-CHECK Propósito:
contrato operativo de health\_check.py --- script de arranque del
sistema, invocado vía alias start. Naturaleza: lectura estricta por
defecto. Única excepción autorizada a escritura: auto-sync condicional
del Entity Index (ver abajo). Ninguna otra sección del script escribe en
Notion, git, ni archivos locales. Checks ejecutados (orden fijo):
version → env → git → vgit → notion → docs\_sync → vdoc → index\_age →
pending\_tickets. \#\#\# KERNEL:HEALTH-CHECK-001 --- Entity Index
Auto-Sync - Umbral: INDEX\_STALE\_THRESHOLD\_HOURS = 24. - Archivos
monitoreados: graph\_v2.json, entity\_index\_v2.json (INDEX\_FILES, en
SCRIPTS\_DIR). - Condición de disparo: al menos un archivo supera el
umbral. - Acción: subprocess a python3 vantage.py sync, cwd =
SCRIPTS\_DIR, timeout 120s. - Frecuencia: máximo una vez por corrida de
health\_check.py, solo si se cruzó el umbral --- no re-sincroniza
índices ya frescos. - Clasificación: housekeeping de rutina, NO
remediación de fallo --- ver KERNEL:FAIL-PHILOSOPHY. Un índice stale no
es un fallo del sistema; es mantenimiento esperado de una estructura de
lectura. - Manejo de error: si el sync falla (returncode ≠ 0, timeout, o
vantage.py no encontrado), el check reporta fail() y el script NO
reintenta ni repara --- a partir de ahí aplica KERNEL:FAIL-PHILOSOPHY
estándar (reportar, esperar instrucción). Justificación de la excepción:
las Golden Rules de ``no reparar automáticamente''
(KERNEL:CV-GOLDEN-RULES) aplican a discrepancias de negocio en el
pipeline (Score, Gate\_Decision, URLs, JD). El Entity Index es
infraestructura de lectura del Runtime, no un dato de negocio --- su
staleness no es un ``fallo'' en el sentido del contrato de Fallos, es
equivalente en naturaleza al sync automático ya existente de L4 (git,
vía launchd) y L3 (Gmail, vía launchd): mantenimiento programado, no
decisión que requiera al operador. \#\#\# KERNEL:HEALTH-CHECK-002 ---
Reporte de Tickets Agrupación por campo Prioridad
(CRÍTICO/ALTO/MEDIO/BAJO/Sin Prioridad) sobre Bug Tracker y Task
Tracker. Detalle explícito (título) solo para CRÍTICO y ALTO;
MEDIO/BAJO/Sin Prioridad solo cuentan. Clasificación Reactivo→Bug /
Proactivo→Task ya definida en KERNEL:TRACKER-SCHEMA. --- \#\#
KERNEL:OWNERSHIP \#\#\# KERNEL:OWNERSHIP-001 --- AI Component -
Responsabilidades del componente de IA (ej: validación de triggers,
generación de HANDOFF). - No modifica campos Class B. \#\#\#
KERNEL:OWNERSHIP-002 --- Python Component - Responsabilidades del
componente Python (ej: escritura en Notion, procesamiento de FEED). -
Único componente con permiso de escritura en Notion. --- \#\#
KERNEL:TRIGGERS 1. TRIGGERS Cada trigger define un contrato de input,
proceso y output. El componente AI no ejecuta pasos fuera del contrato
del trigger activo. \#\#\# KERNEL:TRIGGER-001 FEED Procesamiento por
Lotes. FEED con más de 10 vacantes se divide en lotes de 10. El
procesamiento es secuencial con header de lote. feed\_processor.py no
tiene reintento automático por lote --- ante fallo parcial, reportar
estado y esperar instrucción humana. No hay rollback automático de lotes
previos completados. Origen: Changelog v6.2.1 --- promovido a contrato
activo v8.5.1 \#\#\# KERNEL:TRIGGER-002 VL1 Los comandos VL1 son
wrappers de mantenimiento del Tracker. No son triggers del AI Component
--- son comandos Python autónomos. Se documentan aquí para definir sus
contratos de operación y los límites de lo que ejecutan sin intervención
humana. Restricción de arquitectura: Ningún comando VL1 escribe campos
Class B. - VL1 backfill escribe layer, hash y Prioridad --- campos Class
A. - VL1 batch puede modificar Status --- Class A ---únicamente con
--execute. - VL1 batch --- guardia de escritura: La ausencia del flag
--execute hace el comando permanentemente read‑only. El script no debe
usar input() interactivo como mecanismo de protección --- input() falla
en contextos no‑TTY y puede producir escritura no intencionada. El flag
--execute es el único mecanismo válido de autorización para este
comando. \#\#\# KERNEL:TRIGGER-003 QA Checklist Canónico de 6 ítems: QA
valida formato y completitud del CV exportado. QA no evalúa fit,
oportunidad, score, seniority match, con-veniencia de aplicar ni
alineación estratégica con la vacante. El checklist obligatorio contiene
exactamente 6 ítems: 1. Identidad y contacto 1. Estructura de secciones
1. Orden de experiencia 1. Completitud de contenido 1. Integridad visual
y legibilidad 1. Consistencia de exportación Resultado obligatorio de
QA: GO / NO-GO (Checklist): 1. Identidad y contacto --- PASS/FAIL ---
nota breve 1. Estructura de secciones --- PASS/FAIL --- nota breve 1.
Orden de experiencia --- PASS/FAIL --- nota breve 1. Completitud de
contenido --- PASS/FAIL --- nota breve 1. Integridad visual y
legibilidad --- PASS/FAIL --- nota breve 1. Consistencia de exportación
--- PASS/FAIL --- nota breve Si cualquier ítem retorna FAIL, el
resultado final es NO‑GO. \#\#\# KERNEL:TRIGGER-004 DRY RUN - Campos
Permitidos: Op · Empresa · Rol · URL · Source\_Type · Prioridad · Status
- Campos Prohibidos: Visual Signal · Innovation DNA · Score Estimado ·
Gate\_Decision · Decisión CREATE/BLOCKED \#\#\# KERNEL:TRIGGER-005 SYNC
- Formato de Output (≤12 líneas, sin excepción) - SYNC REPORT ---
{[}FECHA{]} - Target: X \textbar{} Postulado: X \textbar{} En proceso: X
\textbar{} Rechazado: X \textbar{} Total: X - NADs OVERDUE: X - LAST
WRITE: {[}timestamp{]} \#\#\# KERNEL:TRIGGER-006 TOP 3 BY SCORE 1. Marca
\textbar{} Rol \textbar{} Score 1. Marca \textbar{} Rol \textbar{} Score
1. Marca \textbar{} Rol \textbar{} Score - Campos Permitidos: role,
brand, url - Formato de Output: JSON con metadatos de la vacante \#\#\#
KERNEL:TRIGGER-007 NEXT ACTION: \textasciitilde/vantage\_pipeline.sh
status Restricción: SYNC reporta estado. No interpreta tendencias. No
recomienda acciones estratégicas. No compara períodos. Datos puros del
estado actual de Notion. - Campos Permitidos: handoff\_id, status -
Formato de Output: Confirmación de estado \#\#\# KERNEL:TRIGGER-008 FEED
Si recibes JSON de vacantes SIN triggers CV‑A · FAST {[}URL{]} ·
CANON‑UPDATE, responde: ``El procesamiento de FEED está migrado a
feed\_processor.py.'' Excepción FAST: array de longitud 1 + trigger
explícito FAST = procesamiento normal por AI Component. - Campos
Permitidos: query, filters - Formato de Output: Lista de entidades
coincidentes \#\#\# KERNEL:TRIGGER-009 STATUS Ejecuta lectura del estado
general. Responde con el estado del sistema actual. No requiere
escritura ni evaluación. --- \#\# KERNEL:GATE-DECISION \#\#\#
KERNEL:GATE-DECISION-001 --- Lógica de Bypass (precede a toda lógica
estándar)
\texttt{plain\ text\ Source\_Type\ ∈\ \{Inbound,\ Referencia,\ Networking\}\ →\ Gate\_Decision:\ CREATE\ (automático)\ →\ Bypasses:\ URL\_GATE\ +\ Score\ threshold\ +\ Visual\ Signal\ detection\ →\ Razón:\ Un\ contacto\ humano\ verificado\ tiene\ mayor\ señal\ que\ cualquier\ algoritmo}
\#\#\# KERNEL:GATE-DECISION-002 --- Lógica Estándar Solo aplica si no
hay Bypass activo (ver KERNEL:GATE-DECISION-001). Orden de evaluación
(secuencial, no paralelo): 1. URL\_GATE --- primer filtro, precede a
cualquier cálculo de fit. Si el link está muerto o inaccesible → Score =
0, Status = Expirada. Sin excepciones. 1. Score (0--100) --- calculado
por Python sobre VM\_Scope, Role\_Class y match de keywords VM en el JD.
1. Gate\_Decision se deriva del Score: - Score ≥ 60 → CREATE
(Ready-to-Apply) - Score 40--59 → Para Revisar (zona gris, no bloqueado,
requiere juicio humano) - Score \textless{} 40 o VM\_Scope = Off-Target
→ BLOCKED / Archivar Nota: estos thresholds no son constantes editables
en esta sección --- viven en profile\_config.yaml (pesos de scoring) y
en el código de gate\_logic. Esta sección documenta el contrato de orden
y las reglas de decisión, no los valores numéricos exactos de scoring
interno (Class B, propiedad de Python). \#\#\# KERNEL:GATE-DECISION-003
--- Resolución de REVIEW\_NEEDED \textgreater{} ⚠️ ALCANCE DE GAP-03: El
guard GAP-03 protege el pipeline Python (feed\_processor.py →
process\_record()). Escritura directa vía MCP (notion-create-pages /
notion-update-page) y flujos HANDOFF → CV-B no tienen guard equivalente
--- esos puntos de entrada pueden escribir campos Class B sin bloqueo.
Estado: gap documentado, pendiente implementación de class\_b\_guard.py
(FX-1 open). Contrato de Desbloqueo: REVIEW\_NEEDED es un estado de
bloqueo parcial --- la entrada existe en Notion con campos Class A
escritos, pero sus campos Class B están congelados hasta que el operador
resuelva el problema que impidió el procesamiento completo. Disparador
de resolución: Status = ``Target'' es el único valor que
layer\_1\_run.py reconoce como señal de que el operador resolvió el
problema y la entrada está lista para ser procesada. Cualquier otro
valor de Status mantiene el bloqueo. Flujo de resolución --- contrato
formal: 1. Operador corrige el campo problemático en Notion (campo
indicado en Notas). 1. Operador cambia Status → Target. 1. Operador
corre \textasciitilde/vantage\_pipeline.sh. 1. layer\_1\_run.py detecta
Status = Target con Gate vacío o REVIEW\_NEEDED y procesa campos Class B
normalmente: URL\_GATE → Score → Gate\_Decision → VM\_Scope →
Role\_Class. Implementación en código (feed\_processor.py): el
comentario de contrato en process\_record() documenta este flujo
explícitamente. Ver también el guard GAP-03 en el mismo archivo. EXPIRED
(gate decision, campo Class B) ≠ Expirada (operational status, campo
Class A). Son campos distintos con lógica de asignación distinta. El
sistema no los fusiona, no los interpreta como equivalentes, no usa uno
para inferir el otro. \textgreater{} Ejemplo: Un registro puede tener
Status = Expirada (Class A, asignado manualmente o por URL\_GATE en el
primer run) con Gate\_Decision aún vacío --- si Python no ha corrido
todavía. Inversamente, un registro puede tener Gate\_Decision = EXPIRED
(Class B, asignado por Python tras ≥2 runs con URL dead) sin que el
operador haya cambiado Status manualmente. Estos dos estados coexisten
sin conflicto. \#\#\# KERNEL:GATE-DECISION-004 --- Por Qué los Gates Son
Deterministas Un gate que puede sobreescribirse manualmente no es un
gate --- es una sugerencia. La confiabilidad del pipeline depende de que
las decisiones de gate sean predecibles y reproducibles. Si el gate
bloquea, el input de búsqueda necesita ajuste --- no el gate. \#\#\#
KERNEL:GATE-DECISION-005 --- Flujo de Recuperación BLOCKED Gate =
BLOCKED no es estado terminal. RT-1 permite corregir campos Class A
(URL, JD, Source) y re-validar con Python. Si el fix produce CREATE, el
patch se escribe en Notion. RT-1 no sobreescribe el gate; corrige el
input para que Python cambie su decisión. \#\#\#
KERNEL:GATE-DECISION-006 --- REJECTED (Post-Aplicación) REJECTED es un
valor Class B derivado de Status = ``Rechazado'' (Class A, asignado por
el operador cuando la empresa rechaza la postulación). Mismo mecanismo
que APPLIED: el operador escribe una señal Class A observable
externamente; Python la traduce a Class B en el siguiente run de
layer\_1\_run.py (evaluate\_rejection\_status(), análogo a
evaluate\_application\_status()). El operador nunca escribe
Gate\_Decision directamente --- esto no es excepción a
KERNEL:GATE-DECISION-004. Registros con Next\_Action ya poblado quedan
protegidos (PROTECCIÓN TOTAL) y no reciben REJECTED retroactivamente sin
limpieza manual del campo. --- \#\# KERNEL:NAMING-CONVENTION ---
Convención de Nombres de Outputs Todo archivo generado por el sistema
para una vacante específica comparte el mismo stem --- solo la extensión
distingue el tipo de documento (.md, .pdf, .fig). Formato del stem:
\texttt{plain\ text\ \{Año\}\_\{Nombre\}\_\{Apellido\}\_\{Marca\_normalizada\}\_\{Vacante\_normalizada\}}
Reglas de normalización (aplican a Marca y Vacante): - Cada espacio
natural del texto original se reemplaza por guión bajo (\_) --- incluye
espacios entre palabras del mismo campo (ej. ``VM Coordinator'' →
VM\_Coordinator). - Sin acentos ni caracteres especiales (é→e, ñ→n, \& →
``y''). - Sin símbolos de puntuación (---, /, :, comas, paréntesis). -
Guión bajo como único separador en todo el stem --- no se mezcla con
CamelCase. Ejemplo: Vacante: Gucci --- VM Coordinator, LATAM (2026) Stem
resultante: 2026\_Mauricio\_Meyran\_Gucci\_VM\_Coordinator\_LATAM
Archivos de la misma vacante:
\texttt{plain\ text\ 2026\_Mauricio\_Meyran\_Gucci\_VM\_Coordinator\_LATAM.md\ \ \ (CV-B,\ Figma\ tags)\ 2026\_Mauricio\_Meyran\_Gucci\_VM\_Coordinator\_LATAM.pdf\ \ (export\ QA)\ 2026\_Mauricio\_Meyran\_Gucci\_VM\_Coordinator\_LATAM.fig\ \ (si\ aplica,\ archivo\ Figma)}
Aplica a: CV-B (.md), export QA (.pdf), archivo Figma (.fig) y cualquier
output futuro relacionado a una vacante específica. El stem se fija en
el momento de generar el primer entregable (CV-B) y se reutiliza sin
variación en todos los outputs derivados subsecuentes de esa misma
vacante --- el naming no es una decisión de KERNEL:CV-PIPELINE
post-generación, es un contrato que antecede al primer archivo escrito.
No aplica a: DRY RUN archivado (convención propia, ver
KERNEL:ARCHITECTURE-L4 --- ``Archivo DRY RUN archivado mensualmente''),
ni a artefactos de sistema (logs, backups, entity\_index). Relación con
CANON:OUTPUT-CONTRACT-001: son contratos distintos y complementarios.
Output Contract gobierna la estructura interna del contenido (slots,
figma\_text\_id, reglas de serialización). Esta sección gobierna el
nombre físico del archivo en disco. Ninguno reemplaza al otro. --- \#\#
KERNEL:CV-GOLDEN-RULES 1. GOLDEN RULES Límites de Ejecución Las Reglas
de Oro son restricciones de arquitectura. No son preferencias de
comportamiento ni guidelines opcionales. Cada violación genera una
respuesta estandarizada de rechazo. El componente AI no negocia, no
busca workarounds, no ejecuta versiones parciales de una operación
rechazada. \#\#\# KERNEL:CV-GOLDEN-RULES-001 1. No Evaluar Fit Antes de
Escribir El componente AI es executor. La evaluación de fit pertenece a
Python (score determinista) y al humano (decisión final de postulación).
Excepción documentada --- CV‑A: El componente AI extrae keywords y gaps
técnicos para optimización de CV. Esto no es evaluación de fit ni juicio
de oportunidad --- es análisis de alineación técnica para producción del
documento. Solicitudes que activan esta regla: ``¿Es buena esta vacante
para mí?'' ``¿Crees que encajo en este rol?'' ``¿Vale la pena aplicar
aquí?''

Respuesta estandarizada: OPERACIÓN RECHAZADA --- Violación Regla de Oro
\#1 Tu solicitud: {[}descripción{]} Razón: El componente AI no evalúa
fit. El score determinista de Python y tu decisión final son los únicos
evaluadores válidos. Alternativa operativa: Escribe la vacante con FEED
o FAST → \textasciitilde/vantage\_pipeline.sh → revisa Score en
Ready-to-Apply ¿Proceder? Escribe SÍ o CANCELAR

\subsubsection{KERNEL:CV-GOLDEN-RULES-002}\label{kernelcv-golden-rules-002}

\begin{enumerate}
\def\labelenumi{\arabic{enumi}.}
\tightlist
\item
  No Calcular ni Estimar Campos Class B Campos protegidos: Score ·
  VM\_Scope · Role\_Class · Match · Gate\_Decision · Next\_Action ·
  Fetch · Fuente · JD\_Quality · Dedup\_Flag Si el JSON entrante incluye
  valores en estos campos, se ignoran. Si el usuario solicita una
  estimación de score o gate, se rechaza. Python recalcula en cada run
  --- ningún valor estimado por el componente AI tiene validez en el
  pipeline. Solicitudes que activan esta regla: ``¿Qué score crees que
  tendría esta vacante?'' ``¿Pasaría el gate esta entrada?'' JSON con
  ``score'': 75 incluido

  Respuesta estandarizada: OPERACIÓN RECHAZADA --- Violación Regla de
  Oro \#2 Tu solicitud: {[}descripción{]} Razón: Score, Gate y campos
  Class B son Python‑only. Un valor estimado contaminaria el pipeline.
  Alternativa operativa: Escribe la entrada →
  \textasciitilde/vantage\_pipeline.sh → Python calcula con lógica
  determinista ¿Proceder? Escribe SÍ o CANCELAR

  \subsubsection{KERNEL:CV-GOLDEN-RULES-003}\label{kernelcv-golden-rules-003}
\item
  No Cuestionar la Calidad de Datos del Usuario El sistema no comenta
  sobre el volumen de resultados. No sugiere ampliar búsqueda. No evalúa
  si el JSON tiene suficientes entradas. La estrategia de búsqueda es
  100 \% responsabilidad humana. Comportamiento con JSON vacío o de bajo
  volumen: JSON procesado: 0 resultados válidos. No hay nada que
  escribir en Notion. SESIÓN COMPLETADA Sin sugerencias. Sin
  recomendaciones de fuentes alternativas. Sin análisis de por qué el
  resultado fue escaso. Distinción de contexto: Si el JSON llega dentro
  de un flujo DRY RUN ya iniciado (el operador aprobó y el array resultó
  en 0 entradas válidas post‑filtro), el comportamiento es idéntico:
  reportar 0, cerrar sesión. No reiniciar el flujo ni solicitar nuevo
  JSON. \#\#\# KERNEL:CV-GOLDEN-RULES-004
\item
  No Delegar Escritura al Usuario El sistema genera y escribe
  directamente en Notion post‑APROBAR\_WRITE. ``Copia esto y pégalo en
  Notion'' viola esta regla. Excepciones válidas y acotadas: Export PDF
  → fuera del alcance de Notion API Upload a Google Drive → fuera del
  alcance de Notion API Fuera de estas dos excepciones, si el sistema
  puede escribir directamente, escribe directamente. \#\#\#
  KERNEL:CV-GOLDEN-RULES-005
\item
  No Interpretar en SYNC SYNC reporta el estado actual de Notion. Datos
  puros. Sin recomendaciones estratégicas, sin análisis de tendencias,
  sin comparaciones entre períodos, sin sugerencias de próximos pasos
  más allá del output estándar del reporte. Solicitudes que activan esta
  regla dentro de SYNC: ``¿Qué fuentes están funcionando mejor?''
  ``¿Debería ajustar mis targets?'' ``¿Cuál es la tendencia de mis
  scores este mes?''

  Respuesta estandarizada: OPERACIÓN RECHAZADA --- Violación Regla de
  Oro \#5 SYNC reporta datos puros. Análisis e interpretaciones fuera
  del alcance de este trigger. Alternativa operativa: Cierra SYNC → abre
  nueva sesión → solicita análisis con los datos del reporte

  Template Universal de Rechazo OPERACIÓN RECHAZADA --- Violación Regla
  de Oro \#{[}N{]} Tu solicitud: {[}descripción exacta{]} Razón: {[}qué
  regla viola y por qué existe la restricción{]} Alternativa operativa:
  {[}pasos concretos para lograr el objetivo dentro del sistema{]}
  ¿Proceder? Escribe SÍ o CANCELAR

  \begin{center}\rule{0.5\linewidth}{0.5pt}\end{center}

  \subsection{KERNEL:CV-PIPELINE}\label{kernelcv-pipeline}
\item
  CV PIPELINE Contratos de Sesión --- Arquitectura de Dos Sesiones
  Obligatorias \#\#\# CV‑A
\end{enumerate}

\begin{itemize}
\tightlist
\item
  Input: URL o JD de la vacante
\item
  Process: AI Component extrae keywords + identifica gaps + determina
  tono de marca
\item
  Output: HANDOFF (5 campos exactos)
\item
  Cierre obligatorio: SESIÓN COMPLETADA → nueva sesión \#\#\# HANDOFF
  Contrato de Transferencia entre Sesiones (JSON)

  \{ ``empresa'': ``\,``,''rol'': ``\,``,''JD\_keywords\_top6'':
  {[}``\,``,''``,''``,''``,''``,''\,''{]}, ``fit\_gaps'':
  {[}``\,``,''\,''{]}, ``tono\_marca'': ``\,``,''idioma'': ``\,'' \}

  Si cualquier campo está ausente, se solicita. El sistema no inventa
  valores para campos faltantes. Un HANDOFF incompleto no avanza a CV‑B.
  Regla de Idioma: CV‑A detecta el idioma del JD de origen (ES o EN) y
  lo declara en el campo idioma del HANDOFF. Si el JD mezcla ambos
  idiomas, prevalece el idioma predominante por volumen de texto. CV‑B
  usa este valor para seleccionar exclusivamente la versión ES o EN de
  cada sección del Career Canon (CANON:PROFILE-001,
  CANON:EXPERIENCE-001, etc.) --- no se mezclan idiomas en un mismo
  CV‑B. Un HANDOFF sin campo idioma no avanza a CV‑B (mismo tratamiento
  que los demás campos obligatorios). Por Qué Son Dos Sesiones Separadas
  CV‑A es análisis estratégico --- qué posicionar y cómo. CV‑B es
  producción --- el documento final. En una sesión única, el contexto de
  análisis contamina la voz del CV. La separación es una restricción de
  calidad, no de conveniencia. Regla de Orden de Experiencia El orden de
  la experiencia profesional en todos los Derived Outputs es siempre
  cronológico descendente (más reciente primero). El orden no se
  modifica por Positioning Mode, relevancia a la vacante ni ninguna otra
  variable.

  Orden canónico obligatorio: C01 → C02 → C03 → C04 → C05

  Regla de Entrega de Markdown con Figma Tags CV‑B entrega el Markdown
  con Figma tags al operador antes de escribir en Notion. El operador
  revisa y autoriza. Solo tras autorización explícita el AI Component
  escribe el bloque \# MARKDOWN CANON ALIGNED como contenido de la
  página de la vacante en el Tracker. El Markdown no se escribe en
  Notion si el operador no ha autorizado explícitamente. \#\#\# CV‑B
\item
  Input: HANDOFF completo de CV‑A + Career Canon activo
  (notion.so/377938be-fc42-8089-93f2-f52dbd2dec6c)
\item
  Validation: AI Component verifica los 5 campos del HANDOFF antes de
  proceder
\item
  Canon check: AI Component valida que empresa, rol canónico, bullets y
  KPIs del F2 sean derivados del Career Canon --- no inventados ni
  contradictorios con él. Cualquier desviación se reporta antes de
  escribir.
\item
  Nota de Fuente de Verdad --- Skeleton vs Tag Registry
\item
  El Skeleton incluido en esta sección define la estructura visual, el
  orden de bloques y la secuencia obligatoria de contenido para CV‑B.
\item
  Los IDs numéricos exactos, tipos de slot, reglas de inyección y
  condición LOCKED/Variable viven en CAREER\_CANON:OUTPUT-CONTRACT §L
  --- Output Contract. \#\#\# Reglas operativas
\item
  KERNEL:CV-PIPELINE Define la arquitectura de ejecución de CV‑B.
\item
  CANON:OUTPUT-CONTRACT-001 Define el contrato exacto de Figma.
\item
  CV‑B debe usar ambos al mismo tiempo.
\item
  El output final debe preservar un bloque \#\#\#\#\#\# figma\_text\_id
  por cada slot del Skeleton/Tag Registry.
\item
  El orden del output debe coincidir con el Skeleton.
\item
  El significado de cada slot debe coincidir con el Tag Registry.
\item
  Si hay discrepancia entre Skeleton y Tag Registry, se debe detener la
  ejecución y solicitar reconciliación antes de producir F2.
\item
  El literal \#\#\#\#\#\# figma\_text\_id no autoriza inventar, omitir,
  fusionar ni dividir slots. Cada ocurrencia representa un slot
  gobernado por el Tag Registry activo. \#\#\# SKELETON‑INJECTION
  MAPPING (L1 LOGIC)
\item
  El componente AI no tiene permiso para ``decidir'' la estructura
  visual. Su única tarea es el mapping de información del Career Canon
  hacia un Skeleton predefinido en CAREER\_CANON:OUTPUT-CONTRACT
\item
  Invarianza Estructural: Cualquier optimización de CV debe ser una
  copia exacta del Skeleton en cuanto a número de headers y IDs,
  sustituyendo únicamente el contenido textual (payload).
\item
  Auditoría de Estructura: Antes de presentar el resultado final,
  validar: COUNT(figma\_text\_id)\_SKELETON ==
  COUNT(figma\_text\_id)\_OUTPUT. Si los números no coinciden, abortar y
  re‑mapear.
\item
  Auditoría de Secuencia: La auditoría de estructura no es suficiente si
  el count es correcto pero el orden está alterado. Antes de presentar
  el resultado final, verificar que los slots de experiencia aparezcan
  en secuencia canónica estricta: C01 → C02 → C03 → C04 → C05. Ninguna
  variable del HANDOFF --- keywords, tono\_marca, fit\_gaps, Positioning
  Mode --- autoriza alterar esta secuencia. Si el orden no coincide,
  abortar y re‑mapear desde el Skeleton.
\item
  Process: AI Component presenta F2 Markdown completo bajo Output
  Contract v1.0.
\item
  Post‑autorización del operador: AI Component escribe el Markdown como
  contenido de la página de la vacante en Notion bajo encabezado \#
  MARKDOWN CANON ALIGNED.
\item
  Output: Markdown con Figma tags en formato .md descargable → entrega a
  operador para Figma.
\item
  Post‑aplicación: Status = Postulado →
  \textasciitilde/vantage\_pipeline.sh → Python marca APPLIED.
  Componente consumido: Figma Sync recibe el F2 Markdown aprobado por el
  operador como input directo de este pipeline. Arquitectura completa,
  stack e invariantes documentados en KERNEL:ARCHITECTURE-L4 --- no se
  duplica aquí. --- \#\# KERNEL:CANON-UPDATE
\end{itemize}

\begin{enumerate}
\def\labelenumi{\arabic{enumi}.}
\tightlist
\item
  CANON‑UPDATE CANON‑UPDATE actualiza el Career Canon activo. No es una
  operación de discovery, scoring, gate decision ni evaluación de fit.
  Su función es mantener la fuente de verdad profesional alineada con
  nueva evidencia, cambios aprobados por el operador o ajustes de
  estructura requeridos por el Output Contract. Input: Descripción
  explícita del cambio solicitado por el operador. Ejemplos válidos:
\end{enumerate}

\begin{itemize}
\tightlist
\item
  ``Actualizar C01 con nuevo bullet sobre campaña NPI.
\item
  ``Agregar nuevo KPI validado para Levi's.''
\item
  ``Ajustar Positioning Mode N2 para roles de Store Design.''
\item
  ``Actualizar el perfil profesional en español e inglés.''
\item
  ``Modificar Tag Registry porque cambió el Skeleton de Figma.''
  Contexto requerido: Para ejecutar CANON‑UPDATE, el AI Component debe
  cargar:
\item
  KERNEL:CV-PIPELINE
\item
  KERNEL:CV-GOLDEN-RULES

  Si el cambio afecta secciones no incluidas en Runtime, como Education,
  Certifications, Major Projects, Derived Outputs Archive o Changelog,
  el AI Component debe solicitar acceso explícito al CAREER\_CANON.md
  original antes de proceder.

  Validación previa: Antes de modificar cualquier contenido, el AI
  Component debe identificar:
\end{itemize}

\begin{enumerate}
\def\labelenumi{\arabic{enumi}.}
\tightlist
\item
  Qué sección o secciones del Career Canon serán afectadas.
\item
  Qué IDs canónicos se impactan.
\item
  Si el cambio requiere versión ES, EN o ambas.
\item
  Si el cambio impacta CV‑A, CV‑B, QA o el Output Contract.
\item
  Si la información proporcionada es suficiente o requiere confirmación
  del operador.
\item
  Si la información es insuficiente, se debe solicitar aclaración. El
  sistema no inventa datos faltantes. El flujo obligatorio de
  CANON‑UPDATE es:
\item
  Recibir descripción del cambio.
\item
  Identificar secciones afectadas.
\item
  Validar contra Career Canon activo.
\item
  Producir un DRY RUN con:
\item
  Esperar autorización explícita del operador.
\item
  Solo tras APROBAR\_WRITE, producir los dos outputs obligatorios.
\item
  Outputs obligatorios: CANON‑UPDATE siempre produce dos outputs
\item
  Página Notion
\item
  Archivo .md Restricciones
\end{enumerate}

\begin{itemize}
\tightlist
\item
  CANON‑UPDATE no evalúa fit.
\item
  CANON‑UPDATE no calcula score.
\item
  CANON‑UPDATE no modifica campos Class B.
\item
  CANON‑UPDATE no inventa KPIs, fechas, certificaciones, marcas, roles
  ni logros.
\item
  CANON‑UPDATE no altera figma\_text\_id sin instrucción explícita del
  operador.
\item
  CANON‑UPDATE preserva versiones ES/EN cuando la sección afectada sea
  bilingüe.
\item
  CANON‑UPDATE preserva el orden cronológico C01 → C02 → C03 → C04 →
  C05. La sesión termina con:

  CANON-UPDATE COMPLETADO
\end{itemize}

\begin{enumerate}
\def\labelenumi{\arabic{enumi}.}
\tightlist
\item
  Secciones actualizadas:
\end{enumerate}

\begin{itemize}
\tightlist
\item
  {[}lista{]}
\end{itemize}

\begin{enumerate}
\def\labelenumi{\arabic{enumi}.}
\tightlist
\item
  IDs impactados:
\end{enumerate}

\begin{itemize}
\tightlist
\item
  {[}lista{]}
\end{itemize}

\begin{enumerate}
\def\labelenumi{\arabic{enumi}.}
\tightlist
\item
  Outputs entregados:
\item
  Página Notion --- listo / escrito post‑APROBAR\_WRITE
\item
  Archivo .md --- entregado

  Compatibilidad downstream:
\end{enumerate}

\begin{itemize}
\item
  CV‑A: PASS/FAIL
\item
  CV‑B: PASS/FAIL
\item ~
  \subsection{QA: PASS/FAIL}\label{qa-passfail}

  \subsection{KERNEL:FAIL-PHILOSOPHY}\label{kernelfail-philosophy}
\end{itemize}

\begin{enumerate}
\def\labelenumi{\arabic{enumi}.}
\tightlist
\item
  FILOSOFÍA DE FALLO Los fallos del sistema son señales de que el
  pipeline funciona correctamente. No son errores a corregir --- son
  outputs esperados de un sistema de filtrado. Un gate que nunca bloquea
  no está filtrando. La presencia de gates BLOCKED, scores en 0 y
  entradas EXPIRED es evidencia de que el sistema aplica sus criterios
  --- no de que el mercado esté seco o el sistema esté roto. Qué Hace el
  Sistema Cuando Falla No intenta reparar outputs del sistema. No
  sugiere workarounds para entradas bloqueadas. No escala urgencia.
  Reporta el estado y espera instrucción humana para el siguiente paso
  dentro del flujo normal del pipeline. Excepción Documentada --- Gate =
  BLOCKED Gate = BLOCKED recuperable vía RT‑1: si el bloqueo es por
  campos Class A corregibles, RT‑1 es el mecanismo de recuperación. El
  componente AI informa esta opción pero no la ejecuta sin instrucción
  explícita. --- \#\# KERNEL:SCOPE \#\#\# Scope y Economía de Contexto
\end{enumerate}

\begin{itemize}
\item
  Acceso a lógica base preferente vía Terminal (lazy\_loader.py).
\item
  MCP autorizado para lectura, DRY RUN y modificación documental del
  Kernel cuando exista instrucción explícita del operador.
\item
  Terminal continúa siendo la ruta recomendada para operaciones masivas,
  auditorías y cambios estructurales. Runtime: L0 (Lectura estricta).
  Cero escritura directa.
\item
  Jerarquía: L1 \textgreater{} L2 \textgreater{} L3. Claude consolida,
  NO extrae.
\item
  FEED: única vía manual de Claude es FAST. Toda ingesta de L1, L2 y L3
  se realiza metódicamente vía Python (layer\_1\_run.py,
  layer\_3\_mail.py, feed\_processor.py). Ante JSON o FEED sin trigger
  FAST explícito: ``El procesamiento de FEED está migrado a Python; usa
  FAST si requieres entrada manual.''
\item ~
  \subsection{Triaje de ejecución: Antes de usar herramientas, aplicar:
  1. Requerimientos, 2. Triaje de costos (A: Terminal, B: MCP, C:
  Upload), 3. Confirmación. Priorizar Opción
  A.}\label{triaje-de-ejecuciuxf3n-antes-de-usar-herramientas-aplicar-1.-requerimientos-2.-triaje-de-costos-a-terminal-b-mcp-c-upload-3.-confirmaciuxf3n.-priorizar-opciuxf3n-a.}

  \subsection{KERNEL:DATA-FLOW}\label{kerneldata-flow}

  \subsubsection{Flujo de Datos y
  Escritura}\label{flujo-de-datos-y-escritura}
\item
  Pipeline: Kernel → DRY RUN → APROBAR\_WRITE → Notion Write.
\item ~
  \subsection{Pre-validación: Cruzar esquema contra KERNEL:SCHEMA antes
  de cualquier
  escritura.}\label{pre-validaciuxf3n-cruzar-esquema-contra-kernelschema-antes-de-cualquier-escritura.}

  \subsection{KERNEL:ROUTING}\label{kernelrouting}

  \subsubsection{Rutas de Carga (MCP)}\label{rutas-de-carga-mcp}

  Para consultar lógica pesada, prioriza Terminal. Alternativamente, MCP
  puede utilizarse cuando:
\item
  El operador lo solicite explícitamente.
\item
  La operación sea documental.
\item
  Se presente DRY RUN previo.
\item
  Exista autorización posterior mediante APROBAR\_WRITE cuando aplique.
  Ruta recomendada: python lazy\_loader.py --page \{KERNEL\_MASTER\}
  --route \{ruta\} Ruta permitida: MCP.
\item
  ruta: KERNEL:SCHEMA
\item
  ruta: KERNEL:OWNERSHIP
\item
  ruta: KERNEL:TRIGGERS
\item
  ruta: KERNEL:CV-PIPELINE
\item
  ruta: KERNEL:GATE-DECISION
\item
  ruta: KERNEL:CV-GOLDEN-RULES
\item ~
  \subsection{ruta:
  KERNEL:FAIL-PHILOSOPHY}\label{ruta-kernelfail-philosophy}

  \subsection{KERNEL:EVOLUTION}\label{kernelevolution}
\end{itemize}

\begin{enumerate}
\def\labelenumi{\arabic{enumi}.}
\tightlist
\item
  EVOLUCIÓN DEL SISTEMA \#\#\# Deuda Técnica Registrada \textbar{} ID
  \textbar{} Descripción \textbar{} Prioridad \textbar{} Estado
  \textbar{} \textbar{} --- \textbar{} --- \textbar{} --- \textbar{} ---
  \textbar{} \textbar{} DT-014 \textbar{} Extract Runtime Identity
  Contract --- encapsular lógica de entity\_prefix en módulo explícito
  con contrato propio. generate\_entity\_id() actualmente carga desde
  resolver\_registry\_v2.json como fix puntual; la deuda residual es de
  encapsulamiento, no de correctitud. Registrado en release v2.4.0.
  \textbar{} MEDIO \textbar{} Abierto \textbar{} Criterios de Cambio El
  sistema reconoce cuándo un cambio es válido y cuándo no. Esta
  distinción protege la estabilidad arquitectónica del pipeline. Cambios
  válidos --- condiciones que justifican modificación: Cambio
  estructural de mercado: nuevos job boards relevantes, cambios en ATS
  de empresas target Cambio en targets: nuevas empresas, nuevas
  exclusiones, ajuste de geografía Ineficiencia probada con datos:
  bottleneck documentado en pipeline runs Violación de boundary entre
  capas: orchestration haciendo intelligence, sistema calculando campos
  Class B de forma sistemática Cambios inválidos --- condiciones que NO
  justifican modificación: Score ``se siente muy estricto'' → el
  algoritmo determinista es intencional, no un bug Ready‑to‑Apply vacío
  → los inputs de búsqueda necesitan ajuste, no el threshold Un dead
  link apareció → comportamiento normal de mercado, no falla de sistema
  Frustración temporal → el sistema funciona; los inputs necesitan
  revisión Comportamiento ante solicitud de cambio inválido: el
  componente AI identifica la condición como cambio inválido, informa al
  operador la razón (usando la lista anterior), y redirige al workflow
  activo equivalente. No ejecuta el cambio, no negocia, no propone
  alternativas fuera del pipeline. Estabilidad de Arquitectura Central
  Los boundaries de capas no colapsan. La filosofía de verificación no
  se negocia. Los contratos de campo Class A/B no se mezclan. Los
  triggers evolucionan; el scoring puede ajustarse; el schema puede
  expandirse. La arquitectura de tres capas, el URL\_GATE como primer
  filtro y la división de ownership entre AI Component y Python son
  invariantes del sistema. Linaje Histórico --- Preservado, No
  Operacional El sistema mantiene registro de lo que fue construido y
  deprecado: GPT Atlas, Grok discovery, SEARCH‑EXEC/SEARCH‑SIGNAL,
  fórmulas de scoring pre‑v5.0, workflows manuales pre‑v6.0. Se
  reconocen como contexto histórico --- no como código activo, no como
  alternativas válidas al pipeline actual. Mezclar realidad operacional
  con linaje histórico en la misma sesión de procesamiento es un error
  de contexto. Si el usuario referencia un componente legacy, el sistema
  lo identifica como tal y redirecciona al workflow activo equivalente.
  \#\# KERNEL:NORM Contrato de Normalización de IDs:
\end{enumerate}

\begin{itemize}
\tightlist
\item
  Esquema: {[}PREFIX{]}:{[}KEY{]} (ej: KERNEL:TRIGGERS).
\item
  Alcance: Todos los documentos fundacionales (MANUAL, CAREER CANON,
  KERNEL, SYSTEM PROMPT).
\item
  Excepciones: IDs de Notion (UUIDs) en metadatos o URLs.
\item
  Gobernanza: Cambios requieren APROBAR\_WRITE + entrada en Changelog.
  §Reglas de Migración. Ejecutable vía AI Component bajo autorización
  explícita del operador. --\textgreater{} Normalización documental de
  IDs legacy hacia el esquema {[}PREFIX{]}:{[}KEY{]}. Ver
  KERNEL:DOC-CONTRACT para contrato completo y listado de 26 ocurrencias
  (DT-015). \#\# KERNEL:CENSUS-SYNC --- Sincronización Obligatoria del
  ID Census El V-ID-CENSUS es un documento derivado --- su fuente de
  verdad son los IDs reales escritos en los documentos fundacionales
  (Kernel, Manual, Career Canon, System Prompt), no al revés. El Census
  no reemplaza esos documentos ni los precede; los audita. Problema que
  resuelve: sin un gate explícito, un cambio de estado de un ID (⚠️ Stub
  → ✅ Ok) o la creación de un ID nuevo puede quedar reflejado en el
  documento fuente pero no en el Census, generando drift silencioso
  entre lo que el sistema documenta y lo que el Census reporta. Regla 1
  --- Gate de cierre de ticket: Ningún ticket en Bug Tracker o Tasks
  Tracker que implique cambio de estado de un ID (Stub→Ok, creación de
  ID nuevo, deprecación de ID existente) se marca Done sin que el Census
  haya sido regenerado y reflejado ese cambio. Si el re-run de
  generate\_census.py no puede ejecutarse en el momento (ej. sin acceso
  a Terminal), el ticket permanece en estado Blocked-Census --- no se da
  por cerrado en falso. Regla 2 --- Alta de IDs nuevos en el spec +
  deeplink automático: generate\_census.py debe operar en dos modos: (a)
  resolución de IDs ya conocidos en CENSUS\_SPEC, y (b) detección de IDs
  presentes en los documentos fuente que NO están en CENSUS\_SPEC (``IDs
  huérfanos''). Todo ID huérfano detectado se reporta explícitamente al
  operador antes de cerrar el ticket asociado --- no se ignora
  silenciosamente. Para todo ID resuelto (conocido u orfano recién
  agregado), el script genera vía API el deeplink correspondiente al
  bloque exacto en Notion --- la navegación desde el Census hacia la
  porción publicada en el TOC del documento fuente debe ser precisa, no
  aproximada al documento completo. Regla 3 --- Disparo atado a
  Changelog: Toda entrada nueva en V-CHANGELOG que documente cierre de
  tickets con cambio de estado de ID debe ir precedida, en la misma
  sesión, por el re-run de Census. El Census se actualiza antes de que
  Changelog registre el batch --- no después, no como tarea suelta.
  Regla 4 --- Presentación automática de DRY RUN de cierre: Ninguna
  sesión que haya involucrado cambios (constructivos, correctivos o
  destructivos) a la documentación o a bases de datos puede cerrarse sin
  que el AI Component presente, en automático y sin esperar solicitud
  del operador, un resumen DRY RUN de todo lo modificado en la sesión
  --- incluyendo estado de Census, Changelog pendiente/escrito, y
  versión. Esta presentación es un reporte de cierre, no un nuevo write:
  no reabre aprobaciones ya otorgadas, solo consolida y expone lo que
  quedó pendiente o completado. Regla 5 --- Chequeo informativo en
  arranque: health\_check.py reporta la antigüedad del V-ID-CENSUS en
  cada corrida (umbral 7 días), como advertencia amarilla si está
  desactualizado. Este chequeo es puramente informativo --- no bloquea
  el arranque de sesión ni auto-ejecuta generate\_census.py (el script
  pega a la API de Notion con rate-limit real, incompatible con el
  contrato de lectura estricta y rápida de health\_check.py). El gate
  real de obligatoriedad sigue viviendo en el cierre de tickets (Regla
  1), no en el arranque. No aplica a: tickets que no modifican estado de
  ningún ID (ej. fixes de redacción, correcciones de trailing space en
  propiedades Notion). --- \#\# KERNEL:DOC-CONTRACT
\end{itemize}

\begin{enumerate}
\def\labelenumi{\arabic{enumi}.}
\tightlist
\item
  CANONICAL DOCUMENT ID CONTRACT (DOC:CLAVE) Este contrato estandariza
  la referencia cruzada entre componentes del sistema y capas
  documentales, eliminando la dependencia de UUIDs en prompts y lógica
  de negocio. \#\#\# Invariantes del Contrato
\end{enumerate}

\begin{itemize}
\tightlist
\item
  Formato Único: {[}PREFIX{]}:{[}KEY{]} (ej. MANUAL:SETUP).
\item
  Prefix Ownership: Cada prefijo mapea a una única página canónica en
  Notion.
\item
  SSOT: resolver\_registry\_v2.json es la autoridad única para resolver
  Prefijos a UUIDs.
\item
  Resolución Determinista: El Resolver (v1.py) garantiza resolución O(1)
  inyectando el ID crudo al componente solicitante. \#\#\# Prefijos
  Autorizados (v8.9.0) \textbar{} Prefijo \textbar{} Documento Destino
  \textbar{} Mapeo Registry \textbar{} \textbar{} \textbar{} ---
  \textbar{} --- \textbar{} --- \textbar{} --- \textbar{} \textbar{}
  KERNEL \textbar{} V \textbar{} KERNEL \textbar{}
  registry{[}``KERNEL''{]} \textbar{} \textbar{} MANUAL \textbar{} V
  \textbar{} MANUAL \textbar{} registry{[}``MANUAL''{]} \textbar{}
  \textbar{} CANON \textbar{} V \textbar{} CAREER CANON \textbar{}
  registry{[}``CANON''{]} \textbar{} \textbar{} TRACKER \textbar{} V
  \textbar{} TRACKER \textbar{} registry{[}``TRACKER''{]} \textbar{}
  \textbar{} SP \textbar{} V \textbar{} SYSTEM PROMPT \textbar{}
  registry{[}``SP''{]} \textbar{} \textbar{} ALIASES \textbar{} V
  \textbar{} ALIASES \textbar{} registry{[}``ALIASES''{]} \textbar{}
  \textbar{} CHANGELOG \textbar{} V \textbar{} CHANGE LOG \textbar{}
  registry{[}``CHANGELOG''{]} \textbar{} \#\#\# Reglas de Migración Toda
  referencia a páginas del sistema que actualmente use UUIDs
  hardcodeados o anclas planas debe migrar a este esquema.
  lazy\_loader.py es el componente encargado de aplicar este contrato en
  tiempo de ejecución. EXCEPCIÓN DE MIGRACIÓN (DT-015): Se autoriza al
  AI Component a ejecutar la normalización documental
  (search-and-replace) vía MCP sobre las 26 ocurrencias identificadas,
  bajo el trigger NORM {[}DOC:CLAVE{]}.
\end{itemize}

\end{document}
